% Part: sets-functions-relations
% Chapter: relations
% Section: operations

\documentclass[../../../include/open-logic-section]{subfiles}

\begin{document}

\olfileid[ar]{sfr}{rel}{ops}
\olsection{عمليات على العلاقات}

قد ينفع أن نغير العلاقات أو نضم بعضها إلى بعض. ففي
\olref[sfr][rel][ord]{prop:stricttopartial} أخذنا \emph{اتحاد} علاقتين،
وهو اتحادهما حين تؤخذان مجموعتين من الأزواج. وفي
\olref[sfr][rel][ord]{prop:partialtostrict} أخذنا الفرق النسبي
بينهما. وها هنا عمليات أخرى على العلاقات.

\begin{defn}\ollabel{relationoperations}
خذ علاقتين ${\OLId{latin-looped}{R}}$ و${\OLId{latin-looped}{S}}$، ومجموعة كيفما كانت ${\OLId{latin-looped}{A}}$.

نسمّي \emph{معكوس} ${\OLId{latin-looped}{R}}$ المجموعة ${\OLId{latin-looped}{R}}^{-{\OLMathNumeral{1}}} = \Setabs{\tuple{{\OLId{latin-ordinary}{y}}, {\OLId{latin-ordinary}{x}}}}{\tuple{{\OLId{latin-ordinary}{x}},
    {\OLId{latin-ordinary}{y}}} \in {\OLId{latin-looped}{R}}}$.

ونسمّي \emph{الضرب النسبي} لـ${\OLId{latin-looped}{R}}$ و${\OLId{latin-looped}{S}}$ المجموعة $({\OLId{latin-looped}{R}} \mid {\OLId{latin-looped}{S}}) =
\{\tuple{{\OLId{latin-ordinary}{x}}, {\OLId{latin-ordinary}{z}}} : \exists {\OLId{latin-ordinary}{y}}(Rxy \land Syz)\}$.

و\emph{تقييد} ${\OLId{latin-looped}{R}}$ على ${\OLId{latin-looped}{A}}$ معرّف بـ$\funrestrictionto{{\OLId{latin-looped}{R}}}{{\OLId{latin-looped}{A}}}= {\OLId{latin-looped}{R}}
\cap {\OLId{latin-looped}{A}}^{\OLMathNumeral{2}}$.

و\emph{تطبيق} ${\OLId{latin-looped}{R}}$ على ${\OLId{latin-looped}{A}}$ معرّف بـ$\funimage{{\OLId{latin-looped}{R}}}{{\OLId{latin-looped}{A}}} = \{{\OLId{latin-ordinary}{y}} :
(\exists {\OLId{latin-ordinary}{x}} \in {\OLId{latin-looped}{A}})Rxy\}$.
\end{defn}

\begin{ex}
خذ علاقة اللاحق ${\OLId{latin-looped}{S}} \subseteq \Int^{\OLMathNumeral{2}}$ على~$\Int$، وهي
${\OLId{latin-looped}{S}} = \Setabs{\tuple{{\OLId{latin-ordinary}{x}}, {\OLId{latin-ordinary}{y}}} \in \Int^{\OLMathNumeral{2}}}{{\OLId{latin-ordinary}{x}} + {\OLMathNumeral{1}} = {\OLId{latin-ordinary}{y}}}$؛ فـ$Sxy$ يكافئ ${\OLId{latin-ordinary}{x}} + {\OLMathNumeral{1}} = {\OLId{latin-ordinary}{y}}$.

فمعكوسها ${\OLId{latin-looped}{S}}^{-{\OLMathNumeral{1}}}$ علاقة السابق على $\Int$، أي
$\Setabs{\tuple{{\OLId{latin-ordinary}{x}},{\OLId{latin-ordinary}{y}}}\in\Int^{\OLMathNumeral{2}}}{{\OLId{latin-ordinary}{x}} -{\OLMathNumeral{1}} ={\OLId{latin-ordinary}{y}}}$.

وأما ${\OLId{latin-looped}{S}}\mid {\OLId{latin-looped}{S}}$ فتعطي
$ \Setabs{\tuple{{\OLId{latin-ordinary}{x}},{\OLId{latin-ordinary}{y}}}\in\Int^{\OLMathNumeral{2}}}{{\OLId{latin-ordinary}{x}} + {\OLMathNumeral{2}} ={\OLId{latin-ordinary}{y}}}$.

وتقييدها $\funrestrictionto{{\OLId{latin-looped}{S}}}{\Nat}$ علاقة اللاحق على~$\Nat$.

وتطبيقها يعطي $\funimage{{\OLId{latin-looped}{S}}}{\{{\OLMathNumeral{1}},{\OLMathNumeral{2}},{\OLMathNumeral{3}}\}}$ مساوية لـ$\{{\OLMathNumeral{2}}, {\OLMathNumeral{3}}, {\OLMathNumeral{4}}\}$.
\end{ex}

\begin{defn}[الإغلاق المتعدي]
خذ العلاقة الثنائية ${\OLId{latin-looped}{R}} \subseteq {\OLId{latin-looped}{A}}^{\OLMathNumeral{2}}$.

نعرّف \emph{الإغلاق المتعدي} لـ~${\OLId{latin-looped}{R}}$ بأنه ${\OLId{latin-looped}{R}}^+ = \bigcup_{{\OLMathNumeral{0}} < {\OLId{latin-ordinary}{n}} \in
\Nat} {\OLId{latin-looped}{R}}^{\OLId{latin-ordinary}{n}}$؛ والقوى فيه معرفة عوديًا: ${\OLId{latin-looped}{R}}^{\OLMathNumeral{1}} = {\OLId{latin-looped}{R}}$ و${\OLId{latin-looped}{R}}^{{\OLId{latin-ordinary}{n}}+{\OLMathNumeral{1}}} = {\OLId{latin-looped}{R}}^{\OLId{latin-ordinary}{n}}
\mid {\OLId{latin-looped}{R}}$.

وأما \emph{الإغلاق الانعكاسي المتعدي} لـ${\OLId{latin-looped}{R}}$ فتعريفه ${\OLId{latin-looped}{R}}^* = {\OLId{latin-looped}{R}}^+ \cup
\Id{{\OLId{latin-looped}{A}}}$.
\end{defn}

\begin{ex}
في علاقة اللاحق ${\OLId{latin-looped}{S}} \subseteq \Int^{\OLMathNumeral{2}}$ يكون ${\OLId{latin-looped}{S}}^{\OLMathNumeral{2}}xy$ مكافئًا لـ${\OLId{latin-ordinary}{x}} + {\OLMathNumeral{2}} =
{\OLId{latin-ordinary}{y}}$، و${\OLId{latin-looped}{S}}^{\OLMathNumeral{3}}xy$ مكافئًا لـ${\OLId{latin-ordinary}{x}} + {\OLMathNumeral{3}} = {\OLId{latin-ordinary}{y}}$، وعلى هذا القياس.
فـ${\OLId{latin-looped}{S}}^+xy$ إذا وفقط إذا كان ${\OLId{latin-ordinary}{x}} + {\OLId{latin-ordinary}{n}} = {\OLId{latin-ordinary}{y}}$ لعدد ${\OLId{latin-ordinary}{n}} \geq {\OLMathNumeral{1}}$.
وحاصل ذلك أن ${\OLId{latin-looped}{S}}^+xy$ يكافئ ${\OLId{latin-ordinary}{x}} < {\OLId{latin-ordinary}{y}}$، وأن ${\OLId{latin-looped}{S}}^*xy$ يكافئ ${\OLId{latin-ordinary}{x}} \le
{\OLId{latin-ordinary}{y}}$.
\end{ex}

\begin{prob}
بيّن أن إغلاق ${\OLId{latin-looped}{R}}$ المتعدي يحقق التعدية حقًا.
\end{prob}

\end{document}
